\begin{frontmatter}

\title{Climorfa: Climate-Sensitive Urban Morphology in Izmir}

\author[1]{Yusuf Emino\u{g}lu\corref{cor1}\fnref{orcidYusuf}}
\ead{yusuf.eminoglu@deu.edu.tr}
\author[1]{Mert Yava\c{s}\fnref{orcidMert}}
\author[1]{Hilmi Evren Erdin\fnref{orcidHilmi}}
\author[2]{\c{S}\i rk\i \fnref{orcidEsra}}
\cortext[cor1]{Corresponding author.}
\fntext[orcidYusuf]{ORCID: \href{https://orcid.org/0009-0005-6000-2934}{0009-0005-6000-2934}}
\fntext[orcidMert]{ORCID: \href{https://orcid.org/0009-0001-0477-2599}{0009-0001-0477-2599}}
\fntext[orcidHilmi]{ORCID: \href{https://orcid.org/0000-0002-3350-8930}{0000-0002-3350-8930}}
\fntext[orcidEsra]{ORCID: \href{https://orcid.org/0000-0001-7914-4821}{0000-0001-7914-4821}}
\address[1]{Department of City and Regional Planning, Dokuz Eyl\"ul University, \.{I}zmir, T\"urkiye}
\address[2]{Department of City and Regional Planning, Siirt University, Siirt, T\"urkiye}

\begin{abstract}
Local Climate Zones (LCZs) provide a transferable climate-form vocabulary, but
locally actionable interpretation requires finer morphological evidence,
transparent scale choices, and spatially independent validation. We assembled
an analysis-ready 250 m grid for the Izmir functional urban region containing
16,506 grid units and 374 fields from a local 5 m surface model,
building and street geometry, Sentinel-2, Landsat 8/9 land-surface temperature
(LST), Dynamic World, ESA WorldCover, canopy height, coastal context, OpenStreetMap
(OSM), and 2024 population demand. Global LCZ labels were treated as weak labels.
The present analysis retained 5,339 grid units from LCZ 3, 6, 8, and
9 at confidence $\geq 0.60$. LightGBM, XGBoost, Random Forest, and Extra Trees
models were evaluated with five stratified 1.25 km spatial-block folds. Additive ablations tested
morphology, vegetation, green-blue context, full proxy context, and network-cost
two-step floating catchment area (2SFCA) green-space access at 400, 800, and
1,200 m. The primary 800 m LightGBM model achieved macro-F1 0.838, compared
with 0.663 for morphology alone; XGBoost, Random Forest, and Extra Trees reached
0.835, 0.829, and 0.819, respectively, at the primary recipe. Adding vegetation
produced the largest early gain, while 800 m 2SFCA added 0.0056 macro-F1 to
LightGBM relative to an otherwise identical full-context model. LightGBM
leave-one-district-out macro-F1 averaged 0.723 across 11 districts (range
0.475--0.819), revealing substantial transfer heterogeneity. Weak LCZ classes
also differed in summer LST ($H=1592.0$, $\epsilon^2=0.298$); NDVI and canopy
volume were negatively associated with LST, whereas NDBI, building coverage,
and road density were positively associated. The 569-location manual audit was
completed; among 96 high-quality audited units in the four-class modeling population,
exact agreement was 8.3\% for both the weak labels and LightGBM, exposing strong
class-space compression. A municipal park inventory provided strong district-level
macro triangulation of OSM supply but not object-level validation. These findings
establish a reproducible, explainable baseline and show why surface-model structure,
vegetation, network accessibility, spatial transfer, and audit evidence must
accompany LCZ interpretation. Raster/figure-ground deep branches remain untrained,
so deep-fusion superiority is not claimed.
\end{abstract}

\begin{keyword}
    Local Climate Zones \sep urban morphology \sep surface model \sep
    land-surface temperature \sep two-step floating catchment area \sep
    spatial cross-validation \sep explainable machine learning \sep Izmir
\end{keyword}

\end{frontmatter}

\begin{figure}[H]
    \centering
    \includegraphics[width=\textwidth]{graphical_abstract.png}
    \caption*{Graphical abstract summarizing the Climorfa evidence workflow,
    from multisource 250 m feature integration through spatial validation,
    climate-response analysis, and audited-label comparison.}
\end{figure}

